\chapter{External Memory}

\tikzset{
  kn/.style={draw,rectangle split,rectangle split horizontal,rectangle split parts=#1,
    inner sep=3pt,minimum height=6mm,font=\small,fill=blue!6},
  blk/.style={draw,minimum width=1.4cm,minimum height=5mm,fill=gray!10}
}

\section{The External Memory Model}

When data is too large for main memory it lives on disk (or SSD). Accessing a
disk is many orders of magnitude slower than accessing RAM, but data is
transferred in large \emph{blocks}, so reading a whole block costs about the
same as reading one word of it.

\begin{definition}[External memory model (EMM)]
There is an internal memory of size $M$ and an unbounded external memory.
Data moves between them in blocks of $B$ consecutive items. The cost of an
algorithm is the number of \emph{block transfers} (reads and writes); computation
in internal memory is free.
\end{definition}

\begin{center}
\begin{tikzpicture}[arr/.style={{Stealth}-{Stealth},thick}]
  \draw[rounded corners,fill=blue!6] (0,0) rectangle (3,2);
  \node[font=\small,align=center] at (1.5,1.35) {internal memory};
  \node[font=\scriptsize] at (1.5,0.85) {size $M$};
  \node[draw,fill=white,minimum width=1cm,minimum height=5mm,font=\scriptsize] at (1.5,0.35) {CPU};
  \draw[rounded corners,fill=gray!6] (5,-0.3) rectangle (11.4,2.3);
  \node[font=\small] at (8.2,2.0) {external memory (disk)};
  \foreach \r in {0,1} \foreach \c in {0,...,3} \node[blk] at (5.9+\c*1.5,1.2-\r*0.8) {};
  \node[font=\scriptsize] at (5.9,1.2) {$B$ items};
  \draw[arr] (3.1,1) -- node[above,font=\scriptsize]{one block} node[below,font=\scriptsize]{per transfer} (4.9,1);
\end{tikzpicture}
\end{center}

The goal: design algorithms whose block transfers are a factor $B$ (or better)
less than their RAM operation counts, by using every item of each block we read.

\begin{example}
Scanning an array of $n$ items costs $\Theta(n/B)$ transfers. Binary search on a
sorted array costs $\Theta(\log(n/B))$ transfers, barely better than $\log n$, because successive probes are in
different blocks. A balanced BST with nodes scattered on disk costs
$\Theta(\log n)$ transfers per search.
\end{example}

\section{External Sorting}

HeapSort and QuickSort access memory all over the place: $\Theta(n \log n)$
transfers in the worst case. MergeSort is naturally sequential, and adapts well.

\paragraph{Merging.} Merging two sorted runs of total length $\ell$ needs three
blocks of internal memory (two inputs, one output) and $\Theta(\ell / B)$
transfers.

\paragraph{$d$-way merge.} Merge $d$ runs at once: keep one block per run plus an
output block, and a min-heap of the $d$ current front items to pick the next
smallest. Requires $(d + 1)B \le M$, so take $d \approx M/B - 1$.

\begin{lstlisting}
d-way-merge(S1, ..., Sd):          // each Si sorted, read block by block
  P <- min-heap of (front item of Si, i) for i = 1..d
  S <- empty output (written block by block)
  while P is not empty do
    (x, i) <- P.deleteMin()
    append x to S
    remove x from Si
    if Si is not empty then P.insert((front of Si, i))
  return S
\end{lstlisting}

\begin{center}
\begin{tikzpicture}[run/.style={draw,minimum width=2.4cm,minimum height=5mm,font=\scriptsize,anchor=west},
  buf/.style={draw,fill=blue!10,minimum width=0.8cm,minimum height=5mm,font=\scriptsize},
  arr/.style={-{Stealth},thick}]
  \foreach \i/\n/\y in {0/1/0,1/2/-0.9,2/d/-2.1} {
    \node[run] (r\i) at (0,\y) {run $S_\n$ (on disk)};
    \node[buf] (b\i) at (4.0,\y) {block};
    \draw[arr] (r\i) -- (b\i);
  }
  \node at (1.2,-1.4) {$\vdots$};
  \draw[rounded corners,dashed] (3.3,0.75) rectangle (7.6,-2.5);
  \node[font=\scriptsize] at (5.45,0.5) {internal memory};
  \node[draw,circle,fill=orange!20,font=\scriptsize,align=center,inner sep=1pt] (h) at (5.6,-1.0) {min-\\heap};
  \foreach \i in {0,1,2} \draw[arr] (b\i) -- (h);
  \node[buf,fill=green!15] (o) at (7.0,-1.0) {out};
  \draw[arr] (h) -- (o);
  \node[run,fill=green!8] (S) at (8.2,-1.0) {output run (on disk)};
  \draw[arr] (o) -- (S);
\end{tikzpicture}
\end{center}

\paragraph{External MergeSort.}
\begin{enumerate}
  \item Create initial runs: read $M$ items at a time, sort them internally, write
        them back. This gives $n/M$ sorted runs, costing $\Theta(n/B)$ transfers.
  \item Repeatedly $d$-way merge groups of $d \approx M/B$ runs. Each round reads
        and writes everything once ($\Theta(n/B)$ transfers) and divides the
        number of runs by $d$.
\end{enumerate}
Number of rounds: $\log_{d}(n/M)$. Total transfers:
\[
\Theta\!\left(\frac nB \log_{M/B} \frac nM\right).
\]
This is optimal: any comparison-based external sort needs
$\Omega\big(\frac nB \log_{M/B} \frac nB\big)$ transfers.

\begin{example}
With $n = 10^{10}$, $M = 10^8$, $B = 10^4$: initial runs give $100$ runs; one
merge round with $d = M/B = 10^4 \ge 100$ suffices. So the whole sort reads and
writes the data only twice.
\end{example}

\section{External Dictionaries: B-Trees}

To make search efficient in the EMM we want trees whose nodes fill a block and
have many children, reducing the height from $\log_2 n$ to $\log_B n$.

\subsection{Multiway Search Trees}

A node of a \emph{multiway search tree} stores $k$ sorted keys
$k_1 < \cdots < k_k$ and has $k + 1$ subtrees $T_0, \ldots, T_k$, where all keys
in $T_i$ lie between $k_i$ and $k_{i+1}$ (with $k_0 = -\infty$, $k_{k+1} = +\infty$).
Search: at each node find the subtree to descend into (inside the node this is
free in the EMM, or $\bigO(\log k)$ by binary search).

\subsection{2-3 Trees}

\begin{definition}
A \emph{2-3 tree} is a multiway search tree in which
\begin{itemize}
  \item every internal node has either 1 key and 2 children, or 2 keys and 3
        children;
  \item all empty subtrees (leaves) are at the same depth.
\end{itemize}
\end{definition}

\paragraph{Insert.} Search for $k$; it ends at the lowest internal level. Add $k$
to that node. If the node now has 3 keys (\emph{overflow}), \emph{split} it: the
middle key moves up into the parent, and the left and right keys become two
separate nodes. The parent may overflow in turn; if the root splits, a new root
is created and the tree grows by one level (the only way height increases).

\begin{example}
Insert 10, 20, 30 into an empty 2-3 tree: [10, 20] then [10, 20, 30] overflows;
split: root [20] with children [10] and [30]. Insert 40, 50: [30, 40] then
[30, 40, 50] overflows; 40 moves up: root [20, 40] with children [10], [30], [50].
\begin{center}
\begin{tikzpicture}[level distance=11mm,level 1/.style={sibling distance=16mm}]
  \node[kn=1] {20}
    child {node[kn=1] {10}}
    child {node[kn=3,fill=red!12] {30\nodepart{two}40\nodepart{three}50}};
  \node[font=\scriptsize,red!70!black] at (1.2,-1.75) {overflow};
  \draw[-{Stealth},thick] (2.3,-0.6) -- node[above,font=\scriptsize]{split} (3.6,-0.6);
  \begin{scope}[xshift=5.6cm,level 1/.style={sibling distance=16mm}]
  \node[kn=2] {20\nodepart{two}\textcolor{orange!80!black}{\textbf{40}}}
    child {node[kn=1] {10}}
    child {node[kn=1] {30}}
    child {node[kn=1] {50}};
  \end{scope}
\end{tikzpicture}
\end{center}
\end{example}

\paragraph{Delete.} If $k$ is in an internal node, swap it with its successor (in
a bottom node) first. Remove $k$ from the bottom node. If the node becomes
empty (\emph{underflow}):
\begin{itemize}
  \item if an immediate sibling has 2 keys, \emph{transfer} (rotate): a key moves
        from the sibling up to the parent and the parent's separating key moves
        down;
  \item otherwise \emph{merge} with a sibling, pulling the separating key down
        from the parent; the parent may underflow in turn. If the root becomes
        empty, delete it (height decreases).
\end{itemize}

\subsection{(a,b)-Trees}

\begin{definition}
For $2 \le a \le \ceil{b/2}$, an \emph{$(a,b)$-tree} is a multiway search tree
in which
\begin{itemize}
  \item every non-root internal node has between $a$ and $b$ children (so
        between $a-1$ and $b-1$ keys);
  \item the root has between 2 and $b$ children;
  \item all leaves (empty subtrees) are at the same depth.
\end{itemize}
A 2-3 tree is a $(2,3)$-tree; a $(2,4)$-tree is also called a 2-3-4 tree.
\end{definition}

The condition $a \le \ceil{b/2}$ makes splitting work: an overflowing node has
$b$ keys; splitting it around the median yields two nodes with at least
$\ceil{b/2} - 1 \ge a - 1$ keys each.

\begin{theorem}
An $(a,b)$-tree with $n$ keys has height $\Theta(\log_a n)$; more precisely
$h \in \bigO(\log_a n)$ and $h \in \Omega(\log_b n)$.
\end{theorem}
\begin{proof}
Count nodes at each level: the root has at least 2 children, and each other
internal node at least $a$. So level $i \ge 1$ has at least $2a^{i-1}$ nodes, and
the number of leaves (empty subtrees) is at least $2a^{h-1}$. A multiway tree
with $n$ keys has exactly $n + 1$ empty subtrees, so $n + 1 \ge 2a^{h-1}$, giving
$h \le 1 + \log_a\frac{n+1}{2}$. Similarly $n + 1 \le b^{h}$ gives the lower bound.
\end{proof}

Insert and delete work exactly as for 2-3 trees (split on overflow, transfer or
merge on underflow), touching $\bigO(h)$ nodes.

\subsection{B-Trees}

\begin{definition}
A \emph{B-tree of order $m$} is an $(\ceil{m/2}, m)$-tree. In the EMM we choose
$m \in \Theta(B)$ so that one node fills one block.
\end{definition}

\begin{center}
\begin{tabular}{lll}
\toprule
Operation & Block transfers & Internal work\\
\midrule
\code{search} & $\Theta(\log_B n)$ & $\bigO(\log B)$ per node (binary search)\\
\code{insert} & $\Theta(\log_B n)$ & $\bigO(B)$ per node\\
\code{delete} & $\Theta(\log_B n)$ & $\bigO(B)$ per node\\
\bottomrule
\end{tabular}
\end{center}
Since $\log_B n = \frac{\log n}{\log B}$, a B-tree with $B = 1000$ and
$n = 10^{9}$ has height about 3, and the root (and often level 1) can be kept in
internal memory.

\paragraph{Variants.}
\begin{itemize}
  \item \textbf{B$^+$-trees}: all KVPs stored in leaves; internal nodes hold only
        keys for routing, so they have higher fanout; leaves are linked for fast
        range queries. Used by databases and file systems.
  \item \textbf{Pre-emptive splitting/merging}: split full nodes on the way down
        during insert, so a single top-down pass
        suffices.
  \item \textbf{Red-black trees} are a binary encoding of 2-3-4 trees: a 2-3-4
        node becomes a black node with 0, 1 or 2 red children.
\end{itemize}

\section{Extendible Hashing}

Hashing in external memory: we want $\bigO(1)$ block transfers per operation,
and to grow without rehashing everything.

\paragraph{Structure.}
\begin{itemize}
  \item Each key's hash $h(k)$ is a long bit string (say $L$ bits).
  \item A \emph{directory} (kept in internal memory) is an array of size $2^d$,
        where $d$ is the \emph{global depth}. Entry $i$ points to the block for keys
        whose hash starts with the $d$-bit string $i$.
  \item Each block (on disk) stores up to $S \approx B$ keys and has a \emph{local
        depth} $k \le d$: all keys in the block agree on their first $k$ bits.
        A block with local depth $k$ is pointed to by $2^{d - k}$ directory entries.
\end{itemize}

\paragraph{Search.} Read the first $d$ bits of $h(k)$, look up the directory
(free), load one block: 1 transfer.

\paragraph{Insert.} Search for the block $P$. If $P$ has room, add $k$. Otherwise
\emph{split} $P$:
\begin{enumerate}
  \item If $P$'s local depth $k_P < d$: create a new block $P'$, set both local
        depths to $k_P + 1$, redistribute the keys of $P$ by bit $k_P + 1$, and
        redirect half of the directory entries that pointed to $P$ to $P'$.
  \item If $k_P = d$: first \emph{double the directory} ($d \leftarrow d + 1$,
        each old entry is copied into two), then split as above.
\end{enumerate}
Repeat if all keys went to one side.

\begin{lstlisting}
extendible-insert(k):
  P <- block for the first d bits of h(k)
  while P is full do
    if local-depth(P) = d then double the directory; d <- d + 1
    split P into P and P' by bit local-depth(P) + 1
    update directory pointers
    P <- block for the first d bits of h(k)
  insert k into P
\end{lstlisting}

\paragraph{Delete.} Remove $k$ from its block. Optionally merge a block with its
``buddy'' (the block differing only in the last of the $k$ bits) if their
combined contents fit, and halve the directory when all local depths are less
than $d$.

\begin{example}
$S = 2$ keys per block, $d = 1$: block $A$ (prefix 0, $k_A = 1$) holds keys with
hashes \texttt{0010}, \texttt{0111}; block $B$ (prefix 1) holds \texttt{1000}.
Insert \texttt{0101}: $A$ is full and $k_A = d$, so double the directory to
$d = 2$ (entries 00, 01, 10, 11; both 10 and 11 point to $B$, whose local depth
stays 1). Split $A$ by the second bit: $A$ = \{\texttt{0010}\} (prefix 00),
$A'$ = \{\texttt{0111}, \texttt{0101}\} (prefix 01), both with local depth 2.
\end{example}

\begin{center}
\begin{tikzpicture}[de/.style={draw,minimum width=1cm,minimum height=6mm,font=\small\ttfamily},
  bk/.style={draw,rounded corners,fill=blue!6,font=\scriptsize,align=left,minimum width=2.6cm,inner sep=3pt},
  arr/.style={-{Stealth},thick}]
  \node[font=\scriptsize] at (0,0.75) {directory, $d = 2$};
  \foreach \b [count=\i from 0] in {00,01,10,11} \node[de] (e\i) at (0,-\i*0.6) {\b};
  \node[bk] (A) at (4,0.5) {$A$ (local depth 2)\\\texttt{0010}};
  \node[bk,fill=green!10] (A2) at (4,-0.75) {$A'$ (local depth 2)\\\texttt{0111}, \texttt{0101}};
  \node[bk] (B) at (4,-2.0) {$B$ (local depth 1)\\\texttt{1000}};
  \draw[arr] (e0.east) -- (A.west);
  \draw[arr] (e1.east) -- (A2.west);
  \draw[arr] (e2.east) -- (B.west);
  \draw[arr] (e3.east) -- (B.west);
\end{tikzpicture}
\end{center}
After the insertion in the example: $B$ has local depth $1 < d$, so two directory
entries share it.

\paragraph{Analysis.} Search and delete: 1 block transfer. Insert: usually 1-2,
plus occasionally a directory doubling, which is done in internal memory (no
block transfers for the data). No global rehash ever moves all the data; a split
touches only one block. Space: blocks may be partly empty (about 69\% full on
average for random hashes); the directory can blow up if many hashes share long
prefixes.

\section{Summary}
\begin{itemize}
  \item In the EMM we count block transfers of size $B$, with internal memory $M$.
  \item External MergeSort: initial runs of size $M$, then $(M/B)$-way merges:
        $\Theta(\frac nB \log_{M/B} \frac nM)$ transfers.
  \item 2-3 trees, $(a,b)$-trees and B-trees keep all leaves at the same depth
        using split/transfer/merge; a B-tree of order $\Theta(B)$ uses
        $\Theta(\log_B n)$ transfers per operation.
  \item Extendible hashing: a directory on the first $d$ hash bits, blocks with
        local depths; $\bigO(1)$ transfers per operation and local splits instead
        of global rehashing.
\end{itemize}
